% =========================
% chuong/01-gioi-thieu.tex  —  người viết: Trúc
% MỤC 1. GIỚI THIỆU
% Số liệu mô tả phải được tái lập từ pipeline version hóa trước khi nộp báo cáo.
% Cohort: delivered, order_delivered_customer_date <= order_estimated_delivery_date.
% Nhãn: review_score 1--2; chính sách multiple-review hiện giữ câu trả lời mới nhất
% và phải được chốt trong luận văn. Không dùng số liệu từ lần đếm theo chính sách khác.
% =========================

\section{GIỚI THIỆU}

% ------------------------------------------------------------
\subsection{Bối cảnh và lý do thực hiện dự án}

Trong thương mại điện tử, điểm đánh giá của khách hàng là một trong những thông tin
người mua xem trước khi ra quyết định. Nghiên cứu trên hơn 15 triệu đơn hàng của nền
tảng Tmall cho thấy trải nghiệm giao hàng ảnh hưởng đến điểm đánh giá, và điểm đánh giá
đó có liên hệ với khả năng mua hàng của những khách đến sau và doanh số người bán
\parencite{deshpande2023}. Vì vậy, review thấp là một tín hiệu dịch vụ đáng theo dõi,
nhưng dữ liệu quan sát không đủ để quy trực tiếp thành doanh thu bị mất.

Giao hàng trễ là nguyên nhân thường được nhắc đến nhất. \textcite{ravula2023} cho thấy
khách hàng chấm điểm thấp hơn rõ rệt khi đơn đến trễ, trong khi giao sớm chỉ được cộng
thêm rất ít. Do đó, các doanh nghiệp thường theo dõi tỉ lệ giao đúng hạn như một chỉ số
chính của chất lượng dịch vụ. Tuy vậy, giao nhanh hơn chưa chắc làm khách hài lòng hơn:
khi nhà bán lẻ tách đơn để giao nhanh từng phần, sự phiền toái khi nhận nhiều lần giao
có thể triệt tiêu lợi ích của tốc độ và làm giảm việc mua lại \parencite{brylla2022}.

Tuy nhiên, khi phân tích bộ dữ liệu công khai của Olist \parencite{olist2018}, nhóm nhận
thấy vẫn tồn tại các đơn được giao trước hoặc đúng ngày dự kiến nhưng nhận review 1 hoặc
2 sao. Tỷ lệ và số lượng chính thức phải được tái lập từ pipeline theo đúng chính sách
multiple-review đã version hóa trước khi đưa vào báo cáo kết quả. Quan sát này cho thấy
``đúng hạn'' không đồng nghĩa với ``hài lòng'' và tạo ra một nhóm case mà quy tắc chỉ theo
dõi giao trễ sẽ bỏ sót.

Nhóm gọi hiện tượng này là \textit{sự không hài lòng ẩn} (hidden dissatisfaction): khách
hàng không hài lòng dù đơn hàng không có dấu hiệu bất thường nào trên các chỉ số vận hành.
Đây là lý do nhóm chọn đề tài xây dựng một hệ thống web giúp doanh nghiệp nhận diện và xử
lý sớm những đơn hàng này.

% ------------------------------------------------------------
\subsection{Vấn đề cần giải quyết}

Các nghiên cứu dự đoán sự hài lòng trên dữ liệu Olist chủ yếu tập trung so sánh độ chính
xác của các mô hình và dừng lại ở việc gán nhãn cho từng đơn hàng \parencite{orman2026}.
Kết quả dạng này khó được sử dụng trực tiếp trong vận hành, vì nhân viên chăm sóc khách
hàng cần biết đơn nào cần liên hệ trước, vì sao đơn đó có nguy cơ, và nên làm gì tiếp theo.
Ngay cả khi có điểm dự đoán, việc chọn đơn nào để xử lý trước khi năng lực có hạn cũng
không đơn giản: một nghiên cứu gần đây trên ba bộ dữ liệu, trong đó có Olist, cho thấy
xếp hạng bằng mô hình học máy không phải lúc nào cũng tốt hơn quy tắc đơn giản là ưu tiên
đơn có giá trị cao nhất \parencite{li2026}. Vì vậy, người dùng cần một công cụ hiển thị
đồng thời mức rủi ro và các thông tin nghiệp vụ để tự đưa ra quyết định.

Thời điểm can thiệp cũng quan trọng. Trong ngành hàng không, việc chủ động thông báo cho
hành khách trước khi sự cố xảy ra giúp giảm đáng kể phản hồi tiêu cực trên mạng so với
xử lý khi sự cố đã diễn ra \parencite{nazifi2021}. Tương tự, nếu doanh nghiệp liên hệ
khách hàng trước khi họ để lại đánh giá, vấn đề có cơ hội được giải quyết sớm hơn.

Từ góc độ hệ thống thông tin, nhóm xác định ba vấn đề cụ thể:

\begin{itemize}
    \item Chưa có công cụ tập hợp các đơn giao đúng hạn nhưng có nguy cơ bị đánh giá thấp
    thành một danh sách ưu tiên để nhân viên xử lý trước khi khách hàng đánh giá.
    \item Kết quả dự đoán, nếu chỉ ở dạng điểm số trong một tệp dữ liệu, không kèm lý do và
    hành động đề xuất, nên người không chuyên về dữ liệu khó sử dụng.
    \item Không có nơi ghi nhận ai đã xử lý đơn nào và kết quả ra sao, nên doanh nghiệp không
    theo dõi và đánh giá được hiệu quả của việc chăm sóc khách hàng.
\end{itemize}

Từ đó, câu hỏi đặt ra cho dự án là: \textit{cần xây dựng hệ thống web như thế nào để kết quả
dự đoán rủi ro được chuyển thành danh sách ưu tiên, lời giải thích và hành động cụ thể mà
nhân viên chăm sóc khách hàng có thể sử dụng hằng ngày?}

% ------------------------------------------------------------
\subsection{Mục tiêu của dự án}

\textbf{Mục tiêu tổng quát:} xây dựng một hệ thống web hỗ trợ nhân viên chăm sóc khách hàng
phát hiện và xử lý các đơn hàng giao đúng hạn nhưng có nguy cơ nhận đánh giá tiêu cực.

\textbf{Mục tiêu cụ thể:}

\begin{enumerate}
    \item Phân tích và thiết kế hệ thống thông qua các sơ đồ Use Case, DFD, BPMN và ERD.
    \item Thiết kế cơ sở dữ liệu dựa trên 9 bảng dữ liệu Olist, bổ sung các bảng do hệ thống
    tạo ra để lưu kết quả dự đoán, hành động xử lý và người dùng.
    \item Xây dựng bốn màn hình chính: danh sách đơn cần ưu tiên (Risk Queue), chi tiết lý do
    rủi ro của từng đơn (Order Detail), ghi nhận và theo dõi xử lý (Action Tracking), và
    minh bạch nguồn dữ liệu (Data Quality).
    \item Tích hợp output ML cấp product item và kết quả tổng hợp cấp order theo \texttt{max-v1},
    kèm model version, threshold version, aggregation rule và các yếu tố liên quan.
    \item Đảm bảo giao diện hiển thị được trên máy tính, máy tính bảng và điện thoại.
\end{enumerate}

% ------------------------------------------------------------
\subsection{Phạm vi và giới hạn của dự án}

\textbf{Phạm vi:}

\begin{itemize}
    \item Dữ liệu: bộ dữ liệu công khai Brazilian E-Commerce Public Dataset by Olist, gồm
    99.441 đơn hàng từ tháng 9/2016 đến tháng 10/2018, chia thành 9 bảng có quan hệ với nhau.
    \item Đối tượng phân tích: các đơn đã giao đến khách và được giao trước hoặc đúng ngày dự
    kiến. Đơn giao trễ không thuộc phạm vi vì nguyên nhân không hài lòng đã rõ.
    \item Người dùng: nhân viên chăm sóc khách hàng và người quản lý của doanh nghiệp.
    \item Sản phẩm: hệ thống web ở mức nguyên mẫu (prototype), chạy được với dữ liệu thật
    của Olist và kết quả dự đoán rủi ro.
\end{itemize}

\textbf{Giới hạn:}

\begin{itemize}
    \item Báo cáo web không tự huấn luyện lại mô hình; nó nhận artifact từ pipeline ML của
    dự án và tập trung vào contract, trình bày, sử dụng và lưu trữ kết quả.
    \item Olist chỉ có review cấp order, không có nhãn riêng cho product item. Mỗi item kế
    thừa nhãn order; score item chỉ phản ánh rủi ro của order trong bối cảnh item và không
    chứng minh item gây ra review thấp.
    \item Chỉ một phần khách hàng chủ động viết đánh giá, nên những người đánh giá không đại
    diện đầy đủ cho toàn bộ người mua \parencite{li2008}. Đơn không có đánh giá không được
    xem là đơn khách hài lòng.
    \item Khách hàng có rủi ro cao nhất chưa chắc là nhóm phản hồi tốt nhất với việc chăm sóc
    \parencite{ascarza2018}. Vì vậy, hành động đề xuất trên hệ thống chỉ mang tính tham khảo,
    quyết định cuối cùng do nhân viên đưa ra, và dự án không đánh giá hiệu quả thực tế của
    việc liên hệ khách hàng.
    \item Dữ liệu thuộc thị trường Brazil giai đoạn 2016--2018. Kết quả không được xem là áp
    dụng trực tiếp cho thị trường Việt Nam hay thời điểm hiện tại.
    \item Hệ thống chỉ ghi nhận hành động xử lý, không tích hợp gọi điện hay gửi email thật
    đến khách hàng.
\end{itemize}

% ------------------------------------------------------------
\subsection{Phương pháp nghiên cứu và cách tiếp cận}

Nhóm tiếp cận dự án theo hướng phát triển hệ thống thông tin hỗ trợ ra quyết định, gồm
các bước sau:

\begin{enumerate}
    \item \textbf{Khảo sát dữ liệu:} tìm hiểu cấu trúc 9 bảng Olist và đo các số liệu nền
    như tỉ lệ đơn giao đúng hạn bị đánh giá thấp để xác định vấn đề.
    \item \textbf{Phân tích yêu cầu:} xác định người dùng, nhu cầu thông tin và các chức năng
    cần có, thể hiện qua sơ đồ Use Case, DFD và BPMN.
    \item \textbf{Thiết kế:} thiết kế cơ sở dữ liệu (ERD), thống nhất cấu trúc dữ liệu kết
    quả dự đoán, và thiết kế giao diện bằng wireframe trên Figma trước khi lập trình.
    \item \textbf{Xây dựng và tích hợp:} xây dựng giao diện bằng HTML, CSS và JavaScript, đọc
    dữ liệu ở định dạng JSON, kết nối với cơ sở dữ liệu và nạp kết quả dự đoán rủi ro.
    \item \textbf{Kiểm thử và đánh giá:} kiểm tra các chức năng chính và hiển thị trên ba kích
    thước màn hình, sau đó đối chiếu kết quả với các mục tiêu đã đặt ra.
\end{enumerate}

Chi tiết về dữ liệu, công nghệ sử dụng và cách đánh giá được trình bày ở mục 3.
